\chapter{Enquadramento e estado da arte} \label{capenquadramento}

Perante o cenário histórico de vulnerabilidades críticas como as anteriormente mencionadas, a necessidade de ferramentas de agregação torna-se evidente, existindo atualmente várias soluções amplamente utilizadas na indústria. Contudo, estas plataformas distinguem-se do \textbf{SecDash} pelos seus objetivos e complexidade; focam-se habitualmente no segmento \textit{enterprise} e na gestão operacional de ciclo completo, incluindo a orquestração de múltiplos \textit{scanners} e fluxos de remediação (remediation workflows). \\

O SecDash, em contrapartida, não pretende ser uma ferramenta que efetua os seus próprios scans de segurança, mas sim focar-se especificamente na agregação, normalização e visualização centralizada de vulnerabilidades provenientes de múltiplos repositórios em plataformas distintas. O objetivo principal do projeto consiste em fornecer uma integração totalmente funcional com o GitHub e o GitLab, as duas plataformas mais amplamente utilizadas para hospedar código, uma vez que estas abrangem a vasta maioria dos casos de uso em ambientes reais. A integração com pipelines de Jenkins é considerada um \textit{stretch goal}, a ser explorado caso a calendarização assim o permita. \\

\subsection{DefectDojo}
O \textbf{DefectDojo} é uma plataforma DevSecOps que atua como um agregador automático para o seu conjunto de ferramentas de segurança, permitindo organizar facilmente o trabalho de segurança e reportar a postura de segurança da organização aos \textit{stakeholders}. \cite{aboutdefectdojo}\\

De entre as funcionalidades do DefectDojo contam-se, entre outras:
\begin{itemize}
	\item O rastreio e reporte de \textit{Findings} de segurança;
	\item A imposição de SLAs (\textit{Service Level Agreement}) em contexto;
	\item A gestão de \textit{risk-acceptance} e outras decisões de triagem
	\item Coordenação com gestão de testes de penetração tradicionais;
\end{itemize}

Para além de possibilitar a integração com mais de 200 ferramentas de segurança, a sua instalação local requer múltiplos contentores para vários serviços, requerindo inclusive instâncias dedicadas de \textit{Celery} e \textit{Reddis}enquanto \textit{task-queue} para a gestão de tarefas assíncronas tais como o processamento de ficheiros de reporte volumosos ou a sincronização com o \textit{Jira} e envio de notificações aos utilizadores. Estamos portanto perante uma infraestrutura comparativamente mais complexa que a requirida pelo SecDash, que será mais apropriada para equipas pequenas que não requiram todas as funcionalidades disponíveis no DefetDojo.\\

Por outro lado, o modelo de dados utilizado pelo \textbf{DefectDojo} utiliza cinco classes distintas (\textbf{\textit{Product Types}}, \textbf{\textit{Products}}, \textbf{\textit{Engagements}}, \textbf{\textit{Tests}} e \textbf{\textit{Findings}}) para organizar o trabalho das equipas que o utilizem.\cite{modeldefectdojo} Um Engagement, por exemplo, representa um momento no tempo em que um \textit{Test} é realizado, e contem um ou mais \textit{Tests}. A necessidade de instanciar um \textit{Engagement} para cada nova análise introduz uma barreira de agilidade significativa. Apesar de se apresentar como flexível, de forma a que este modelo se possa adaptar às diferentes equipas que o utilizem, a abordagem do DefectDojo revela uma preocupação mais voltada para questões de auditoria do que para a visão imediata de vulnerabilidades e agilidade durante o processo de desenvolvimento do software.\\

\subsection{ArmorCode}
O \textbf{ArmorCode} é uma aplicação comercial \textit{enterprise} de \textit{Application Security Posture Management} (ASPM). que oferece integração com mais de 350 ferramentas de segurança de aplicações, infraestrutura, cloud e contentores, bem como com sistemas de DevOps. \\

De entre as suas vastas funcionalidades faz também parte a agregação de vulnerabilidades encontradas por outras ferramentas de \textit{scanning}. De forma semelhante ao SecDash, o ArmorCode não faz ele próprio um scan de vulnerabilidades, apresentando-se como uma solução «\textit{scannerless, vendor-agnostic}»\cite{scanerless}. Mais uma vez, no entanto, estamos perante uma ferramenta de elevada complexidade, possuindo imensas funcionalidades para lá do escopo do problema que o SecDash pretende resolver e que para além disso podem levar a uma curva de aprendizagem acentuada que poderá ser contraproducente quando se procura uma solução ágil para a simples agregação de relatórios de vulnerabilidades externos. \\

Quando comparado com o DefectDojo temos ainda a agravante de esta ser uma ferramenta SaaS proprietária com um modelo de licenciamento orientado para grandes corporações, o que a torna inacessível para projetos independentes ou de equipas de pequena dimensão. \\
